\documentclass[conference]{IEEEtran}
\IEEEoverridecommandlockouts

\usepackage[T1]{fontenc}
\usepackage[utf8]{inputenc}
\usepackage{amsmath}
\usepackage{booktabs}
\usepackage{array}
\usepackage{graphicx}
\usepackage{subcaption}
\usepackage{listings}
\usepackage{xcolor}
\usepackage{url}
\usepackage{hyperref}
\hypersetup{hidelinks}

\definecolor{lstbg}{gray}{0.96}
\definecolor{lstkw}{rgb}{0.10,0.10,0.55}
\definecolor{lstcm}{rgb}{0.30,0.45,0.30}
\lstset{
  basicstyle=\ttfamily\scriptsize,
  backgroundcolor=\color{lstbg},
  keywordstyle=\color{lstkw}\bfseries,
  commentstyle=\color{lstcm}\itshape,
  breaklines=true,
  breakatwhitespace=true,
  columns=fullflexible,
  keepspaces=true,
  frame=single,
  framesep=3pt,
  rulecolor=\color{gray!40},
  xleftmargin=2pt,
  showstringspaces=false,
  postbreak=\mbox{\textcolor{gray}{$\hookrightarrow$}\space},
}
\lstdefinestyle{sh}{language=bash}

\newcommand{\param}[1]{\texttt{#1}}

\begin{document}

\title{Reproducible Processing of Home Blood-Pressure Time Series:\\
Methodology and Operation of the Generator \texttt{generate\_bp\_tikz.py}}

\author{%
\IEEEauthorblockN{Erwin Grabler}
\IEEEauthorblockA{ARTCOM Vienna, Vienna, Austria\\
\texttt{OE1EKG}}}

\maketitle

\begin{abstract}
Home blood-pressure monitoring (HBPM) typically produces an irregularly sampled
time series: several readings on some days, none on others. A naive average over
all individual readings therefore over-weights measurement-intensive days and
distorts the overall picture. This document describes the methodology and
operation of \texttt{generate\_bp\_tikz.py}, a command-line tool that turns a
simple CSV file into two coordinated LaTeX/PGFPlots diagrams and
machine-readable statistics CSVs. The core principle is a two-stage
aggregation: first by calendar day, then optionally into fixed-length blocks.
We explain the computation variants (calendar-day-weighted mean versus median of
daily medians), the outlier logic, and the generated data formats, and we
document every parameter in tabular form. The goal is a reproducible,
transparent, and clinically interpretable representation that explicitly shows
comparison thresholds without implying aneurysm- or patient-specific target
values.
\end{abstract}

\begin{IEEEkeywords}
Blood pressure, HBPM, time series, aggregation, median, interquartile range,
LaTeX, PGFPlots, reproducibility.
\end{IEEEkeywords}

\section{Introduction and Motivation}
Self-measurement of blood pressure at home yields an irregularly sampled time
series spanning weeks or months. Two methodological problems recur. First, the
number of readings per day varies strongly; a day with eight readings would be
weighted eight times as heavily as a day with a single reading in a naive
overall average. Second, real HBPM data contain individual strongly deviating
values (for example immediately after physical exertion) that affect robust and
non-robust location measures differently.

\texttt{generate\_bp\_tikz.py} addresses both points through a deterministic,
parameterizable pipeline. The input is a CSV file; the output consists of a
fragment that can be included in an existing LaTeX document, two standalone
documents, and two optional statistics CSVs. All location and dispersion
measures are computed explicitly and traceably in the source, so the figures can
be checked directly against the underlying numbers.

\section{Data Model and Input}
A CSV file is expected with at least the columns \emph{Date}, \emph{Systolic},
and \emph{Diastolic}; optional columns \emph{Pulse}, \emph{Time}, and
\emph{Note} are recognized and evaluated where useful. Column detection uses
alias names in German and English (e.g.\ \texttt{date}/\texttt{datum},
\texttt{systolic}/\texttt{systolisch}), with case and special characters
normalized.

Dates are accepted in the formats \texttt{YYYY-MM-DD}, \texttt{DD.MM.YYYY},
\texttt{DD.MM.YY}, \texttt{DD/MM/YYYY}, and \texttt{DD/MM/YY}. Numeric values are
parsed in both English and European notation: values such as \texttt{130},
\texttt{130.5}, \texttt{130,5}, \texttt{1.234,56}, and \texttt{1,234.56} are
interpreted correctly by treating the last separator that occurs as the decimal
separator. Non-numeric suffixes such as ``mmHg'' are tolerated. The input
delimiter is auto-detected by default (comma, semicolon, or tab) and can be
forced with \param{-{}-delimiter}.

On input, only rows within the chosen date window (\param{-{}-date-from} through
optional \param{-{}-date-to}) are accepted. Missing or invalid systolic or
diastolic values cause a clear error abort with the row number, rather than
silently producing corrupted statistics.

\section{Aggregation Methodology}

\subsection{Stage 1: Aggregation per Calendar Day}
In the first step, all readings of a calendar day are combined into a daily
statistic. For systolic and diastolic pressure as well as for the derived pulse
pressure (the difference between systolic and diastolic), the minimum, maximum,
median, and arithmetic mean are determined; analogously for pulse, if present.
This daily reduction is the decisive step against over-weighting
measurement-intensive days: regardless of whether a day contains one or eight
readings, it subsequently enters the further computations as \emph{one} daily
value.

The daily diagram (Figure~1 of the output) shows the daily median as a point and
the daily range (minimum to maximum) as a vertical bar for each day. Days with
only one reading therefore have no visible range. The median is preferred over
the mean here because, within a day, it is robust against a single outlier and,
together with the displayed range, captures both the typical level and the
dispersion of the day.

\subsection{Stage 2: Block Aggregation}
In the second step, the daily values are grouped into blocks of fixed length
(default seven days, configurable via \param{-{}-block-days}). Block boundaries
are counted from \param{-{}-date-from}, so the first block starts independently
of the weekday. The last block may be shorter; the caption explicitly notes when
it comprises fewer days and is therefore less robust.

For each block, two distinct quantities are computed and deliberately displayed
separately in the diagram:
\begin{itemize}
\item the \emph{central line} as the block's location measure, and
\item the \emph{interquartile range (IQR)} of the daily medians as a dispersion
measure, drawn as an error bar from $Q_1$ to $Q_3$.
\end{itemize}
The quartiles are determined by linear interpolation of the sorted daily
medians, comparable to the default of common statistical packages.

\section{Computation Variants of the Central Line}
The central methodological decision concerns the location measure of the block
central line. The tool offers two variants via \param{-{}-week-central}, each
serving a clear purpose.

\subsection{Mean (Default)}
In the default (\texttt{mean}), the central line is the calendar-day-weighted
mean of the daily means. Formally, for a block with $D$ days and daily means
$\bar{x}_d$:
\begin{equation}
\mu_{\mathrm{block}} = \frac{1}{D}\sum_{d=1}^{D} \bar{x}_d .
\end{equation}
Since $\bar{x}_d$ is already the mean of the readings of day $d$,
$\mu_{\mathrm{block}}$ is a two-stage, day-equal-weighted mean. This variant is
compatible with the clinically common HBPM/ESC mean thresholds, which are also
based on means, and is therefore suitable for direct comparison with such
thresholds.

\subsection{Median of Daily Medians}
With \texttt{median}, the central line is the median of the block's daily
medians:
\begin{equation}
m_{\mathrm{block}} = \operatorname{median}\bigl(\{\tilde{x}_d\}_{d=1}^{D}\bigr),
\end{equation}
where $\tilde{x}_d$ is the median of day $d$. This variant is more robust
against individual outlier days and is also consistent with the displayed IQR
box, since both are based on the same daily medians. It is recommended when the
typical level should be shown despite occasional exertion spikes.

\subsection{Recommendation}
The choice is not a matter of right or wrong but of purpose: for comparison with
mean-based thresholds, \texttt{mean} is natural; for an outlier-robust trend
display, \texttt{median}. Because dispersion (IQR box) and location are shown
separately, the information remains fully visible in both cases.

\section{Outlier Marking}
Optionally (\param{-{}-week-outliers}), the weekly diagram marks individual days
whose daily median exceeds or falls below a clinical threshold as small points
next to the respective IQR marker. The upper thresholds default to 135\,mmHg
systolic and 85\,mmHg diastolic, matching the HBPM comparison lines drawn in the
diagram; they are adjustable via \param{-{}-week-outlier-sys-hi} and
\param{-{}-week-outlier-dia-hi}. Additionally, optional \emph{lower} thresholds
can be activated via \param{-{}-week-outlier-sys-lo} and
\param{-{}-week-outlier-dia-lo}, for example to reveal days that are too low
under therapy. This marking does not replace a statistical outlier definition
but highlights clinically interesting individual days without removing them from
the aggregation; all days remain included in mean, median, and IQR.

\section{Reference Elements and Thresholds}
Both diagrams contain the same reference elements, explained once in a shared
paragraph and referenced by the captions. The light-gray corridors
120--129\,mmHg systolic and 70--79\,mmHg diastolic are general ESC orientations
under therapy subject to individual tolerability and are explicitly \emph{not}
aneurysm-specific target lines. The dotted lines at 135\,mmHg systolic and
85\,mmHg diastolic mark frequently used home comparison thresholds. This
deliberate separation between comparison threshold and individual therapy target
prevents the figure from implying an unsupported, patient-specific target line.

The automatic summary line at the bottom of the daily diagram states the
day-weighted mean, the median of daily medians, and the number of days with a
daily mean below 135/85\,mmHg. It can be suppressed with
\param{-{}-no-daily-summary-label}.

\section{Generated Files}
By default, the tool writes an includable LaTeX fragment (\param{-{}-out}), a
single-page standalone document with both diagrams stacked
(\param{-{}-standalone-out}), and a ``two-sides'' standalone document in which
each diagram occupies its own tightly cropped page (\param{-{}-two-sides-out}).
Optionally, two statistics CSVs are written (\param{-{}-daily-stats},
\param{-{}-weekly-stats}). The single-page document scales the diagram heights at
typesetting time from the available page height, so the page is always fully
used regardless of the length of the captions.

Tables~\ref{tab:daily_cols} and~\ref{tab:weekly_cols} document the columns of the
two statistics CSVs.

\begin{table}[!t]
\caption{Columns of the daily-statistics CSV (\texttt{-{}-daily-stats}).}
\label{tab:daily_cols}
\centering
\footnotesize
\renewcommand{\arraystretch}{1.15}
\begin{tabular}{>{\raggedright\arraybackslash}p{0.30\columnwidth} >{\raggedright\arraybackslash}p{0.58\columnwidth}}
\toprule
\textbf{Column} & \textbf{Meaning} \\
\midrule
\texttt{date} & calendar day (ISO 8601) \\
\texttt{day\_index} & day index from \texttt{date-from} (0-based) \\
\texttt{n} & number of readings on this day \\
\texttt{sys\_min/max} & systolic minimum/maximum \\
\texttt{sys\_median} & systolic daily median \\
\texttt{sys\_mean} & systolic daily mean \\
\texttt{dia\_min/max} & diastolic minimum/maximum \\
\texttt{dia\_median} & diastolic daily median \\
\texttt{dia\_mean} & diastolic daily mean \\
\texttt{pulse\_*} & pulse min/max/median/mean (if present) \\
\texttt{pulse\_pressure\_*} & pulse pressure min/max/median/mean \\
\bottomrule
\end{tabular}
\end{table}

\begin{table}[!t]
\caption{Columns of the block-statistics CSV (\texttt{-{}-weekly-stats}).}
\label{tab:weekly_cols}
\centering
\footnotesize
\renewcommand{\arraystretch}{1.15}
\begin{tabular}{>{\raggedright\arraybackslash}p{0.40\columnwidth} >{\raggedright\arraybackslash}p{0.48\columnwidth}}
\toprule
\textbf{Column} & \textbf{Meaning} \\
\midrule
\texttt{block} & block index (0-based) \\
\texttt{start\_date}, \texttt{end\_date} & block start/end \\
\texttt{x\_mid} & x position (day index) of block center \\
\texttt{n\_days} & number of days in the block \\
\texttt{n\_readings} & number of readings in the block \\
\texttt{label} & block axis label \\
\texttt{sys\_mean\_of\_daily\_means} & weighted systolic mean \\
\texttt{sys\_q1/q3\_daily\_medians} & systolic $Q_1$/$Q_3$ of daily medians \\
\texttt{dia\_mean\_of\_daily\_means} & weighted diastolic mean \\
\texttt{dia\_q1/q3\_daily\_medians} & diastolic $Q_1$/$Q_3$ of daily medians \\
\texttt{pulse\_pressure\_*} & analogous for pulse pressure \\
\bottomrule
\end{tabular}
\end{table}

\section{Anonymized Application Example}
Figures~\ref{fig:daily} and~\ref{fig:weekly} show the two generated diagrams on
an anonymized, realistic data set (101 readings on 37 calendar days). The figure
labels are produced by the tool in German; an English gloss is given in the
captions. Figure~\ref{fig:daily} is the daily course with daily median, daily
range, the automatic summary line, and the \texttt{n=} labels (readings per
day). Figure~\ref{fig:weekly} is the condensed 7-day overview; here the median
central line (\param{-{}-week-central median}) and outlier marking
(\param{-{}-week-outliers}) were chosen. In both figures the ESC corridors and
HBPM comparison lines are drawn as a shared reference.

\begin{figure*}[!t]
\centering
\includegraphics[width=0.92\textwidth]{fig1.pdf}
\caption{Anonymized example, daily course. Points show the daily median,
vertical bars the daily range (min--max); gray \texttt{n=} labels at the top
give the readings per day. Summary line glosses: \emph{Tage seit}~=~days since;
\emph{Mittelwert}~=~mean; \emph{Median}; \emph{Tage}~=~days;
\emph{Blutdruck}~=~blood pressure.}
\label{fig:daily}
\end{figure*}

\begin{figure*}[!t]
\centering
\includegraphics[width=0.92\textwidth]{fig2.pdf}
\caption{Anonymized example, 7-day overview with median central line and outlier
marking. Lines show the median of daily medians, error bars the interquartile
range of daily medians; circles mark individual days whose daily median exceeds
the threshold (135/85\,mmHg). Legend glosses: \emph{IQR der Tagesmediane}~=~IQR
of daily medians; \emph{7-Tage-Median}~=~7-day median;
\emph{Ausrei{\ss}er-Tage}~=~outlier days; \emph{HBPM-Vergleich}~=~HBPM
comparison.}
\label{fig:weekly}
\end{figure*}

Figure~\ref{fig:onepage} additionally shows the single-page standalone document
(\param{-{}-standalone-out}), in which both diagrams, the shared reference
paragraph, and the captions are combined on one page. The diagram heights are
computed at typesetting time from the available page height, so the page is fully
used regardless of the caption length.

\begin{figure}[!t]
\centering
\fbox{\includegraphics[width=0.97\columnwidth]{fig3.pdf}}
\caption{Anonymized example of the single-page standalone document
(\texttt{-{}-standalone-out}): both diagrams, the shared reference paragraph, and
the captions on one page, with the diagrams scaled to the page height. The figure
labels are produced by the tool in German; see Figures~\ref{fig:daily}
and~\ref{fig:weekly} for English glosses of the recurring terms.}
\label{fig:onepage}
\end{figure}

Finally, Figure~\ref{fig:twosides} shows the two-sides variant
(\param{-{}-two-sides-out}), in which each diagram occupies its own tightly
cropped page and the corresponding caption is set as a text node below the axis.
This format is particularly suited to embedding the individual diagrams as
separate appendix figures.

\begin{figure*}[!t]
\centering
\begin{subfigure}{0.48\textwidth}
\centering
\fbox{\includegraphics[width=\textwidth]{fig4a.pdf}}
\caption{Page 1: daily course.}
\label{fig:twosides_a}
\end{subfigure}
\hfill
\begin{subfigure}{0.48\textwidth}
\centering
\fbox{\includegraphics[width=\textwidth]{fig4b.pdf}}
\caption{Page 2: 7-day overview.}
\label{fig:twosides_b}
\end{subfigure}
\caption{Anonymized example of the two-sides variant
(\texttt{-{}-two-sides-out}): each diagram occupies its own cropped page with the
caption set as a node below the axis. The figure labels are produced by the tool
in German; see Figures~\ref{fig:daily} and~\ref{fig:weekly} for English glosses.}
\label{fig:twosides}
\end{figure*}

\section{Parameter Reference}
Table~\ref{tab:params} lists all command-line parameters with default and
purpose.

\begin{table*}[!t]
\caption{Complete parameter reference of \texttt{generate\_bp\_tikz.py}.}
\label{tab:params}
\centering
\footnotesize
\renewcommand{\arraystretch}{1.25}
\begin{tabular}{>{\raggedright\arraybackslash}p{0.24\textwidth} >{\raggedright\arraybackslash}p{0.16\textwidth} >{\raggedright\arraybackslash}p{0.52\textwidth}}
\toprule
\textbf{Parameter} & \textbf{Default} & \textbf{Purpose} \\
\midrule
\param{-{}-csv} & --- (required) & Input CSV with date, systolic, diastolic (optional time/pulse/note). \\
\param{-{}-date-from} & --- (required) & First date in the evaluation window. \\
\param{-{}-date-to} & all from \texttt{date-from} & Last date; if omitted, all subsequent values. \\
\param{-{}-out} & \texttt{bp\_diagrams.tex} & Path of the includable LaTeX fragment. \\
\param{-{}-daily-stats} & off & Optional daily statistics as CSV. \\
\param{-{}-weekly-stats} & off & Optional block statistics as CSV. \\
\param{-{}-block-days} & \texttt{7} & Block length of the second aggregation stage, in days. \\
\param{-{}-delimiter} & \texttt{auto} & Delimiter: \texttt{auto}/\texttt{comma}/\texttt{semicolon}/\texttt{tab}. \\
\param{-{}-standalone-out} & \texttt{bp\_\hspace{0pt}diagrams\_\hspace{0pt}both\_\hspace{0pt}onepage\_\hspace{0pt}standalone.tex} & Single-page standalone with both diagrams. \\
\param{-{}-no-standalone} & off & Do not write the single-page standalone. \\
\param{-{}-standalone-title} & \texttt{Blutdruckdiagramme} & Title above the standalone diagrams. \\
\param{-{}-show-daily-n} & off & Gray \texttt{n=} labels (readings per day) in the daily diagram. \\
\param{-{}-daily-n-y} & automatic & Fixed y position of the \texttt{n=} labels. \\
\param{-{}-week-outliers} & off & Mark outlier days in the weekly diagram. \\
\param{-{}-week-outlier-sys-hi} & \texttt{135.0} & Upper systolic outlier threshold (mmHg). \\
\param{-{}-week-outlier-dia-hi} & \texttt{85.0} & Upper diastolic outlier threshold (mmHg). \\
\param{-{}-week-outlier-sys-lo} & off & Optional lower systolic threshold. \\
\param{-{}-week-outlier-dia-lo} & off & Optional lower diastolic threshold. \\
\param{-{}-week-central} & \texttt{mean} & Central line: \texttt{mean} (weighted mean) or \texttt{median} (median of daily medians). \\
\param{-{}-no-daily-summary-label} & off & Suppress the automatic summary line in the daily diagram. \\
\param{-{}-two-sides-out} & \texttt{bp\_\hspace{0pt}diagrams\_\hspace{0pt}standalone\_\hspace{0pt}two\_\hspace{0pt}sides.tex} & Two-sides standalone (one cropped page per diagram). \\
\param{-{}-no-two-sides} & off & Do not write the two-sides standalone. \\
\param{-{}-two-sides-width-cm} & \texttt{16.0} & Fixed axis width (cm) of the two-sides variant. \\
\bottomrule
\end{tabular}
\end{table*}

\section{Operation}
\subsection{Prerequisites}
The including LaTeX document requires:
\begin{lstlisting}[style=sh]
\usepackage{tikz}
\usepackage{pgfplots}
\pgfplotsset{compat=1.18}
\end{lstlisting}
The script itself has no external Python dependencies beyond the standard
library.

\subsection{Minimal Call}
\begin{lstlisting}[style=sh]
python3 generate_bp_tikz.py \
  --csv bp_raw.csv \
  --date-from 2026-05-15
\end{lstlisting}
Produces the fragment and both standalone documents with default names.

\subsection{With Date Window and Statistics Export}
\begin{lstlisting}[style=sh]
python3 generate_bp_tikz.py \
  --csv bp_raw.csv \
  --date-from 15.05.2026 \
  --date-to 20.06.2026 \
  --out bp_fragment.tex \
  --daily-stats bp_daily_stats.csv \
  --weekly-stats bp_weekly_stats.csv
\end{lstlisting}

\subsection{Median Central Line with Outliers}
\begin{lstlisting}[style=sh]
python3 generate_bp_tikz.py \
  --csv bp_raw.csv \
  --date-from 2026-05-15 \
  --week-central median \
  --week-outliers
\end{lstlisting}

\subsection{Fully Specified Call}
\begin{lstlisting}[style=sh]
python3 generate_bp_tikz.py \
  --csv bp_raw.csv \
  --date-from 2026-05-15 \
  --date-to 2026-06-20 \
  --week-central median --week-outliers \
  --week-outlier-sys-hi 135 \
  --week-outlier-dia-hi 85 \
  --show-daily-n \
  --out bp_fragment.tex \
  --standalone-out bp_onepage.tex \
  --two-sides-out bp_two_sides.tex \
  --two-sides-width-cm 16 \
  --daily-stats bp_daily_stats.csv \
  --weekly-stats bp_weekly_stats.csv
\end{lstlisting}

\subsection{Including the Output}
The fragment is included in running text with \verb|\input{bp_fragment.tex}|.
Alternatively, the single-page or two-sides standalone can be compiled
separately and the resulting PDF embedded as an appendix figure.

\section{Reproducibility and Validation}
Since all location and dispersion measures are computed explicitly in the source
and output via the statistics CSVs, the figures can be checked directly against
the numbers. A typical validation step is to compare the summary line of the
daily diagram (weighted mean, median of daily medians, number of days below
135/85\,mmHg) with the values in the daily-statistics CSV. The same input, the
same date window, and the same parameters lead deterministically to identical
output.

\section{Implementation and Tested Environment}
The tool is implemented in Python~3 and uses only the standard library (no
third-party dependencies); Python~3.9 or newer is recommended. The generated
documents were typeset with \texttt{pdflatex} (TeX~Live) using \texttt{tikz} and
\texttt{pgfplots} (\texttt{compat=1.18}); the accompanying IEEE papers use
\texttt{IEEEtran.cls} version~V1.8b.

Testing was performed \emph{only} on the following platforms: (i)~macOS with
desktop Python and a desktop TeX~Live/MacTeX installation; (ii)~iPad and iPhone
with \emph{a-Shell}, a local iOS terminal that bundles Python~3 and TeX~Live
(with TikZ and LuaTeX) -- both the script and the LaTeX compilation run
on-device; and (iii)~\emph{Texifier} (formerly TeXpad) on macOS and iPadOS/iOS
for editing and typesetting the LaTeX documents. Note that Texifier's built-in
live typesetter (TexpadTeX) runs in a sandbox and does not permit shell-escape
(\texttt{\textbackslash write18}); the Python script itself is therefore run in
a-Shell or an ordinary shell, while Texifier is used to typeset the generated
\texttt{.tex} files. Other platforms (Windows, Linux) are expected to work with a
standard Python~3 and TeX~Live installation but have not been tested.

\section{Limitations}
The tool is a processing and visualization tool, not a diagnostic system. The
thresholds drawn are comparison and orientation lines; the determination of
individual target values rests with the treating physician. The expressiveness
of short blocks with few days is limited and is marked accordingly in the
caption. The outlier marking is a threshold-based highlighting and not a
statistical outlier classification.

\section{Summary}
\texttt{generate\_bp\_tikz.py} converts irregularly sampled HBPM time series via
a two-stage aggregation into two coordinated, reproducible diagrams. The
separate display of location (mean or median) and dispersion (range or IQR), the
transparent outlier logic, and the explicit separation of comparison threshold
and therapy target enable a clinically interpretable yet traceable
representation. Full parameterization via the command line makes the tool
suitable both for one-off evaluations and for recurring, automated reports.

\section*{Acknowledgment and Disclosure of Generative AI Use}
In accordance with the IEEE guidelines on the use of AI-generated content, the
following is disclosed. The AI system Claude (Anthropic) was used in the
preparation of this document. The AI was used in an assistive capacity for: the
linguistic drafting and structuring of the running text; the formulation of the
methodology description based on the source code developed by the author; the
creation of the LaTeX and table structure; and the preparation of the anonymized
example figures. The conception, software, underlying data, technical claims, and
final responsibility for the content rest entirely with the author; all
AI-assisted contributions were reviewed, verified, and corrected where necessary
by the author. The AI is not an author of this work. A more detailed description
of the AI contribution is provided in the file \texttt{AI\_USAGE.md} of the
accompanying repository.

\end{document}
